\section{Introduction}

Simplicial sets have been used extensively in algebraic topology as a way to encode complex information about topological spaces into a \emph{relatively} simple combinatorial object, which often allows for more streamlined study of aspects of the space. More recently, they've become extremely useful in category theory, especially homotopy theory and higher category theory, due to their ability to encode homotopic information. A well known example of this is quasi-categories, a certain type of simplicial set which forms a convenient model for $(\infty ,1)$-categories, which have been studied extensively by Andr\'e Joyal in \cite{joyal1} and \cite{joyal2}, and Jacob Lurie in \cite{htt}.

Unfortunately, simplicial sets are just complicated enough to be difficult for beginners to work through easily, but are just simple enough for experts to consider them fairly trivial objects, and as such there are very few expositions on simplicial sets which are readily understandable by the amateur. For example, the references \cite{gj} and \cite{may} are considered standard on the topic, but the former is quite advanced for a first pass, and the latter is described by this author as being ``the densest textbook in existence.''

Friedman recently wrote \cite{fried} which we believe to be a very good introduction, but it focuses much more on the topological aspects of the theory than we would like. We also like Emily Riehl's \cite{sset} which is very categoricaly motivated, but we feel her introduction strays too far from the geometric roots of the subject for a first introduction. Given the prominence of simplicial sets in category theory, we would like for there to be an entry level resource that focuses on the categorical motivations of the theory, without fully forsaking the increduble intuition lended to the subject by geometry. This paper is our attempt to produce such a resource.

Category theory is the main prerequisite of this paper. Any student with an understanding of categories, limits, adjunctions, and an undergrad understanding of point-set topology should be more than equipped to read this paper. We recommend \cite{ctx} and \cite{lein} as good resources to learn category theory from, and \cite{ml} as an excellent reference, although we discourage learning from here.

Throughout the paper, we work in the standard category $\Top$ whenever possible, and only pass to the category $\C\mathcal{G}$ of compactly generated spaces when needed. We do this to avoid overloading readers who may not have had advanced topology exposure, especially since passing to $\C\mathcal{G}$ is a largely technical point. Most of the theory developed herein works the same in $\Top$ as in $\C\mathcal{G}$ up to homotopy equivalence.
